\chapter{命题即类型与基础证明}

\section{Curry--Howard 对应}

Lean 中命题位于 \code{Prop}，证明命题 $P$ 就是构造一个类型为 $P$ 的项。常见逻辑连接词对应类型构造：
\begin{center}
\begin{tabularx}{0.94\textwidth}{>{\bfseries}l l X}
\toprule
逻辑 & Lean 形式 & 证明数据 \\
\midrule
$P\Rightarrow Q$ & \code{P -> Q} & 把 $P$ 的证明映到 $Q$ 的证明的函数 \\
$P\land Q$ & \code{And P Q} & 同时含 $P$ 与 $Q$ 证明的结构 \\
$P\lor Q$ & \code{Or P Q} & 左分支或右分支，并携带相应证明 \\
$\forall x, P(x)$ & \code{forall x, P x} & 对每个 $x$ 返回证明的依赖函数 \\
$\exists x, P(x)$ & \code{Exists fun x => P x} & 一个 witness 和它满足性质的证明 \\
$\neg P$ & \code{Not P} & 从 $P$ 的证明导出 \code{False} 的函数 \\
\bottomrule
\end{tabularx}
\end{center}

这一观点把“证明步骤”变成精确的构造操作，而不是自然语言中模糊的“显然”。

\section{直接证明与 tactic proof}

同一个定理可以直接写 proof term，也可以用 tactic 逐步构造：
\begin{lstlisting}
example (p q : Prop) (hp : p) (hq : q) : And p q :=
  And.intro hp hq

example (p q : Prop) (hp : p) (hq : q) : And p q := by
  constructor
  exact hp
  exact hq
\end{lstlisting}
第二段执行 \code{constructor} 后产生两个目标。编辑器展示的 proof state 可以抽象成：
\[
\underbrace{p,q:\PropT,\ hp:p,\ hq:q}_{\text{local context}}
\quad\vdash\quad
\underbrace{p\land q}_{\text{target}}.
\]
每个 tactic 都应当被理解为如何改变这个 judgment。

\section{蕴含、全称量词与引入}

证明函数类型时，引入输入：
\begin{lstlisting}
example (p q : Prop) : p -> q -> p := by
  intro hp
  intro hq
  exact hp

example : forall n : Nat, n = n := by
  intro n
  rfl
\end{lstlisting}
\code{intro} 不是“假设结论成立”，而是按照函数类型构造一个接受证明或对象的函数。

\section{合取、析取与分类讨论}

合取的使用是拆字段，析取的使用必须覆盖两个分支：
\begin{lstlisting}
example (p q : Prop) (h : And p q) : And q p := by
  cases h with
  | intro hp hq => exact And.intro hq hp

example (p q : Prop) (h : Or p q) : Or q p := by
  cases h with
  | inl hp => exact Or.inr hp
  | inr hq => exact Or.inl hq
\end{lstlisting}
\code{cases} 按 \code{And.intro} 构造子拆出两个字段。Lean 也支持更紧凑的结构模式，但它不会减少需要提供的数据。

\section{存在量词必须给出 witness}

\begin{lstlisting}
example : Exists fun n : Nat => n + 1 = 4 := by
  refine Exists.intro 3 ?_
  rfl
\end{lstlisting}
\code{refine} 提供带“洞”的部分证明，\code{?\_} 变成新目标。存在性证明的核心是 witness $3$，而不是 tactic 的名字。

\section{否定、矛盾与底类型}

\code{False} 没有构造子；一旦得到它，可以推出任意命题：
\begin{lstlisting}
example (p : Prop) (hp : p) (hnp : Not p) : False := by
  exact hnp hp

example (p q : Prop) (h : p -> q) (hnq : Not q) : Not p := by
  intro hp
  exact hnq (h hp)
\end{lstlisting}
第二条定理的目标 \code{Not p} 展开就是 \code{p -> False}，所以先 \code{intro hp}。

\section{等式、重写与简化}

若有 $h:a=b$，\code{rw [h]} 按方向把 $a$ 替换为 $b$；\code{rw [<- h]} 反向替换。\code{simp} 使用标记过的简化规则反复重写：
\begin{lstlisting}
example (a b c : Nat) (h : a = b) : a + c = b + c := by
  rw [h]

example (xs : List Nat) : [] ++ xs = xs := by
  simp
\end{lstlisting}
\code{rw} 更适合表达“使用了哪条等式”；\code{simp} 更适合清理单位元、构造子和已注册规则。遇到失败时，先明确期望哪一处发生哪种替换。

\section{常用 tactic 的逻辑角色}

\begin{center}
\begin{tabularx}{0.94\textwidth}{>{\ttfamily}lX}
\toprule
\normalfont tactic & 主要作用 \\
\midrule
intro & 为函数、蕴含或全称目标引入变量/假设 \\
exact & 提供类型与当前目标完全匹配的证明项 \\
apply & 反向使用定理，把当前目标变成它的前提 \\
constructor & 使用目标类型的构造子并产生字段目标 \\
cases & 按归纳类型或析取的构造子分支讨论 \\
rcases & 用模式拆解存在、合取或结构 \\
rw & 按给定等式定向重写 \\
simp & 用简化规则集做规范化重写 \\
have & 声明局部中间结论，控制证明结构 \\
show & 重新表述当前目标，帮助 elaboration 或阅读 \\
\bottomrule
\end{tabularx}
\end{center}

\section{构造逻辑与经典逻辑}

Lean 的核心逻辑是构造性的。排中律、双重否定消去等经典原理可以显式导入或局部使用，但应让依赖可见：
\begin{lstlisting}
example (p : Prop) : Or p (Not p) := by
  classical
  exact Classical.em p
\end{lstlisting}
形式化数学中使用经典逻辑通常完全合理；重要的是不要把“可构造 witness”与“经典地证明存在”混为一谈，尤其当后续希望从证明提取算法时。

\begin{warningbox}
Tactic 成功不等于你已经理解定理。至少应能回答：当前目标是什么、tactic 使用了哪个构造子或定理、产生了哪些子目标、是否引入经典公理，以及删除哪个假设后证明会失败。
\end{warningbox}

\begin{checkpointbox}
\begin{enumerate}
  \item 分别用 proof term 和 tactic 证明合取交换；
  \item 证明 \code{(p -> q) -> (Not q -> Not p)}，逐行解释 proof state；
  \item 构造一个自然数 witness，证明某个简单存在命题；
  \item 用 \code{rw} 与 \code{simp} 各完成一个等式证明，并比较两者；
  \item 找出一个需要经典逻辑的命题，明确标出 \code{classical} 的作用域。
\end{enumerate}
\end{checkpointbox}
